\section{Introduction}

Automated assessment of surgical technical skill is motivated by the prospect of scalable, consistent feedback during training. If such a system is to inform training or credentialling decisions, it must generalise to surgeons who were not seen during its development. Whether current evaluation practice on the field's standard benchmark establishes that it does is the question we take up.

The concern is not specific to surgery. When a model is trained and tested on data from the same individuals, it can succeed by recognising the individual rather than the property of interest --- a failure mode named identity confounding \cite{chaibubneto2019detecting} and demonstrated systematically by Saeb et al. \cite{saeb2017need}, whose activity-recognition error stayed near 2\% when records from the same subjects appeared on both sides of the split, against roughly 7\% when whole subjects were held out and 27\% when only two subjects were available. Electroencephalography supplies the closest structural parallel to what we report below: an Alzheimer's classifier reaching 99.8\% accuracy under segment-based holdout fell to 53.0\% on unseen subjects, with a confidence interval spanning chance, and a seizure classifier fell from 79.1\% to 65.1\% \cite{brookshire2024leakage}. The same pattern appears in speech, where depression detection degrades sharply on unseen speakers because the acoustic features are entangled with speaker identity. These are instances of shortcut learning \cite{geirhos2020shortcut}, and leakage of this general kind has been documented across 17 fields, collectively affecting 294 papers \cite{kapoor2023leakage}. We return to the methodological debate this literature provoked in Section~2.3.

The benchmark itself ships two cross-validation protocols: one holds out a repetition while every surgeon remains in training, the other holds out a surgeon entirely. Both are sanctioned and both appear in this literature, so what follows is not a criticism of an informal practice. That performance falls under the second is not in dispute; it has been reported repeatedly. What has not settled is why. Three studies of JIGSAWS attribute the same gap to three different causes: limited generalisation, answered by a call for surgeon-invariant features; the quality of the dataset's representations and annotations; and an absence of subject-independent validation in earlier work. A fourth, on open-surgery hand tracking rather than robot kinematics, comes closest to the account we pursue, describing its own gap as quantifying ``the surgeon-specific component captured by video-level splits'' --- and stops there. Three different attributions, none of them measured, and a fourth study that names a mechanism without measuring it either. The test is inexpensive: train the same features to predict which surgeon performed the trial, and read off how much identity is recoverable. Such classifiers exist --- surgeons have been identified from intraoperative force profiles with high accuracy --- but they have been built to enable per-surgeon feedback, not to ask whether a skill benchmark is measuring skill. To our knowledge the test has not been run as a validity check on a skill-assessment protocol. The gap has been treated as a limitation to engineer away rather than as a quantity to measure.

Judging such a system also requires care about the target it is fitted to. The reference standard is a human rating, and human ratings of technical skill are neither perfectly reliable nor uniformly so: pooled inter-rater reliability for global rating scales in simulation-based assessment has been reported at 0.78 (95\% CI 0.71-0.83, 23 studies) \cite{ilgen2015checklists}, while untrained raters have been measured at 0.26, with the benefit of rater training lost during periods of non-use \cite{gawad2019interrater}. On this benchmark specifically, the self-declared skill labels track the expert ratings poorly \cite{hendricks2024exploring}. A model here is therefore fitted to a noisy and training-dependent judgment rather than to a fixed standard. Two consequences follow for how results should be read. Agreement statistics cannot be weighed against a single fixed bar. And a reported accuracy or correlation is not self-interpreting: it needs an explicit comparison with a predictor that uses no input at all, a point we return to in Section~3.5 with a concrete demonstration of how misleading its absence can be.

This paper asks what a skill model evaluated on JIGSAWS is actually measuring, and whether the confidence such a model reports would support the human-in-the-loop use these systems are proposed for. We evaluate two prediction paths on identical information --- a conventional baseline of 16 kinematic features with ridge and logistic regression, and a commercial zero-shot decision model given the same features rendered as text --- under both protocols the benchmark sanctions. We make four contributions. First, we quantify the protocol gap with paired surgeon-level intervals: changing only the evaluation protocol, with features, model and data held constant, removes at least a third of the reported skill-classification accuracy (Section~4.1). Second, in a post-hoc diagnostic, we show that the same features identify which of eight surgeons performed a held-out trial with accuracy between 0.83 and 0.93 against chance rates near 0.13, which distinguishes identity confounding from the underfitting explanation that a performance gap alone leaves open (Section~4.2). Third, we report what is to our knowledge the first calibration and selective-prediction audit of a surgical skill assessment model built from motion data, finding both paths overconfident on unseen surgeons and neither path's confidence able to identify the trials it gets wrong (Section~4.5, Section~4.6). Fourth, we benchmark both against a single-feature regression on trial duration, which matches them on the one task of three where any cross-surgeon signal survives (Section~4.4). We do not claim to have discovered identity confounding: the mechanism is well established elsewhere, and what we contribute is its demonstration, quantification and consequences in this setting.
